\section{Spur Gears}
\label{sec:spur_gears}

This section assembles the parametric curves derived in the previous sections into a complete spur gear tooth. This includes determining start and end values of curve parameters $\phi$ necetetating deriving curve intersections. Lastly all results are collected in a summary.

\subsection{Tooth Geometry and Positioning of Curves}
\label{sec:tooth_geometry_positioning}
To construct a gear tooth, the two parametric curves (involute and undercut curve) need to be correctly positioned relative to each other and to the $x$-axis. The curves should be rotated such that mirroring them along the $x$-axis produces a valid tooth geometry with both a right and left flank.

\subsubsection{Tooth Thickness and Half Tooth Angle, \texorpdfstring{$\gamma$}{γ}}
\label{subsec:tooth_thickness}
The tooth thickness (in arc length) at the pitch circle, assuming \textbf{no profile shift}, is:
\begin{equation}
    s_0 = \frac{m\pi}{2}
\end{equation}

\begin{figure}[ht]
    \centering
    \includegraphics[width=\textwidth]{de-evolventenverzahnung-zahnrad-profilverschiebung-vergroesserung-zahndicke.jpg}
    \caption{Gemetric description of additional tooth width due to profile shift. Here $V=x\cdot m :=$ Absolute profile shift value. (\textcite{profileshift})}
    \label{fig:profile_shift_geometry}
\end{figure}
The additional tooth width due to profile shift is more easily computed on the cutting rack as the linear addiional width corresponds to the additional arc length of the tooth on the gear. Figure~\ref{fig:profile_shift_geometry} geometrically shows how the additional tooth width, $\Delta s$ is computed:
\begin{equation}
    \Delta s = m\cdot x \cdot \tan(\alpha_n)
\end{equation}
Resulting in:
\begin{subequations}
\begin{align}
    s_0 &= \frac{p_0}{2} + \Delta s \\
        &= \frac{m\pi}{2} + m \cdot x \cdot \tan(\alpha_n) \\
        &= \frac{m}{2}\left(\pi + 2 \cdot x \cdot \tan(\alpha_n)\right)
\end{align}
\end{subequations}
Therefore, the angle spanned by half of the tooth at the pitch circle is:
\begin{subequations}
\begin{align}
    \gamma_0    &= \frac{s_0}{d_p} \\
                &= \frac{m\left(\pi + 2 \cdot x \cdot \tan(\alpha_n)\right)}{2\,d_p}
\end{align}
\end{subequations}

To get the angle spanned by half the tooth at the base circle, the angle spanned by the involute at $\phi(d_p)$ must be added.

The equation for obtaining the involute angle as derived in Section~\ref{subsec:involute_angle} is:
\begin{equation}
\theta(d^*) = \left[\left(\frac{d^*}{d_b}\right)^2 - 1\right]^{\tfrac{1}{2}} - \arctan\left(\left[\left(\frac{d^*}{d_b}\right)^2 - 1\right]^{\tfrac{1}{2}}\right)
\end{equation}

Inserting the relevant diameters:
\begin{subequations}
\begin{align}
\Delta \gamma &= \theta(d_p) - \cancelto{0}{\theta(d_b)} \\
    &= \theta(d_p) \\
    &= \left[\left(\frac{d_p}{d_b}\right)^2 - 1\right]^{\tfrac{1}{2}}
        - \arctan\left(\left[\left(\frac{d_p}{d_b}\right)^2 - 1\right]^{\tfrac{1}{2}}\right)
\end{align}
\end{subequations}

This results in a half tooth angle at the base circle of:

\begin{subequations}
\begin{align}
    \gamma &= \gamma_0+\Delta \gamma \\
           &= \frac{m\left(\pi + 2 \cdot x \cdot \tan(\alpha_n)\right)}{2\,d_p} + \left[\left(\frac{d_p}{d_b}\right)^2 - 1\right]^{\tfrac{1}{2}} - \arctan\left(\left[\left(\frac{d_p}{d_b}\right)^2 - 1\right]^{\tfrac{1}{2}}\right)
\end{align}
\end{subequations}

\subsubsection{Positioned Involute}
The involute simply needs to be rotated by half the tooth angle at the base circle:
\begin{subequations}
\begin{align}
    \vb{r_{inv,pos}}(\phi) &= \vb{R}(\mypm(-\gamma))\vb{r_{inv}}(\phi) \\[1ex]
    &= \begin{bmatrix}
            \cos(\gamma) & \sin(\mypm\gamma) \\
            -\sin(\mypm\gamma) & \cos(\gamma)
        \end{bmatrix}
        \left(
            \tfrac{d_b}{2} \begin{bmatrix}
                \cos(\phi) \\
                \sin(\phi)
            \end{bmatrix}
            + \tfrac{d_b}{2}\phi \begin{bmatrix}
                \sin(\phi) \\
                -\cos(\phi)
            \end{bmatrix}
        \right)
\\[1ex]
    &= \frac{d_b}{2} \begin{bmatrix}
            \cos(\gamma) & \sin(\mypm\gamma) \\
            -\sin(\mypm\gamma) & \cos(\gamma)
        \end{bmatrix}
        \begin{bmatrix}
            \cos(\phi) +\phi \sin(\phi)\\
            \sin(\phi) -\phi \cos(\phi)
        \end{bmatrix}
\\[1ex]
    &= \frac{d_b}{2} \begin{bmatrix}
            \phantom{-}\cos(\myppm\gamma)\;\cos(\phi) +\cos(\myppm\gamma)\;\phi \sin(\phi) + \sin(\mypm\gamma)\;\sin(\phi) - \sin(\mypm\gamma)\;\phi\cos(\phi)\\
            -\sin(\mypm\gamma)\;\cos(\phi) - \sin(\mypm\gamma)\;\phi \sin(\phi) + \cos(\myppm\gamma)\;\sin(\phi) - \cos(\myppm\gamma)\;\phi\cos(\phi)\\
        \end{bmatrix}
\end{align}
\end{subequations}

\subsubsection{Positioned Undercut Curve}

Due to the definition of the undercut curve (see Figure~\ref{fig:undercut_video_screenshot}), it needs to be rotated by both the half tooth angle at the base circle and the transversal pressure angle ($\alpha_t$).

\begin{subequations}
\begin{align}
    &\vb{r_{undercut,pos}}(\phi) = \vb{R}(\mypm(-\gamma - \alpha_t))\;\vb{r_{undercut}}(\phi)
\\[1ex]
    &= \begin{bmatrix}
            \cos(\gamma+\alpha_t) & \sin(\mypm(\gamma+\alpha_t)) \\
            -\sin(\mypm(\gamma+\alpha_t)) & \cos(\gamma+\alpha_t)
        \end{bmatrix}
        \;\frac{1}{2}\left(
            \vb A \cos(\phi) + \vb B \sin(\phi) + d_p\,\phi\;\vb t(\phi)
        \right)
\\[1ex]
    &= \frac{1}{2}\begin{bmatrix}
            \cos(\gamma+\alpha_t) & \sin(\mypm(\gamma+\alpha_t)) \\
            -\sin(\mypm(\gamma+\alpha_t)) & \cos(\gamma+\alpha_t)
        \end{bmatrix}
        \left(
            \begin{bmatrix}
                a \\ b
            \end{bmatrix} \cos(\phi)
            + \begin{bmatrix}
                -b \\ a
            \end{bmatrix} \sin(\phi)
            + d_p\,\phi\; \begin{bmatrix}
                \sin(\phi) \\ -\cos(\phi)
            \end{bmatrix}
        \right)
\\[1ex]
    &= \frac{1}{2} \begin{bmatrix}
            \cos(\gamma+\alpha_t) & \sin(\mypm(\gamma+\alpha_t)) \\
            -\sin(\mypm(\gamma+\alpha_t)) & \cos(\gamma+\alpha_t)
        \end{bmatrix}
        \begin{bmatrix}
            a\cos(\phi) - b \sin(\phi) + d_p\,\phi\sin(\phi) \\
            b\cos(\phi) + a \sin(\phi) - d_p\,\phi\cos(\phi)
        \end{bmatrix}
\\[1ex]
    &= \frac{1}{2}\bigg[\begin{array}{cc}
            \phantom{-}a\cos(\myppm\;\gamma+\alpha_t\;)\;\cos(\phi)
            - b\cos(\myppm\;\gamma+\alpha_t\;)\;\sin(\phi)
                + d_p\cos(\myppm\;\gamma+\alpha_t\;)\;\phi\sin(\phi)\dots \\
            -a\sin(\mypm(\gamma+\alpha_t))\;\cos(\phi)
                +b\sin(\mypm(\gamma+\alpha_t))\;\sin(\phi)
                - d_p\sin(\mypm(\gamma+\alpha_t))\;\phi\sin(\phi)\dots
        \end{array}
\nonumber\\
        &\phantom{= \frac{1}{2}\bigg[}\begin{array}{cc}
            \hspace{1em}\dots + b\sin(\mypm(\gamma+\alpha_t))\;\cos(\phi)
                + a\sin(\mypm(\gamma+\alpha_t))\;\sin(\phi)
                - d_p\sin(\mypm(\gamma+\alpha_t))\;\phi\cos(\phi)\\
            \hspace{1em}\dots+ b\cos(\myppm\;\gamma+\alpha_t\;)\;\cos(\phi)
                + a\cos(\myppm\;\gamma+\alpha_t\;)\;\sin(\phi)
                - d_p\cos(\myppm\;\gamma+\alpha_t\;)\;\phi\cos(\phi)
        \end{array}\bigg]
\end{align}
\end{subequations}

\subsection{Curve Intersections}
\label{sec:curve_intersections}

The Newton-Raphson method used to numerically solve for both intersections is described in Section~\ref{subsec:newton_raphson}.

\subsubsection{Left--Right Involute Intersection}
In certain configurations, particularly when profile shift is applied, the left and right involute flanks of a tooth may intersect before reaching the dedendum circle. When this occurs, the involute curves must be truncated at their intersection point rather than extending to the dedendum. To determine this intersection, the positioned involute equations are set equal.

\paragraph{System of Equations}
The right and left involutes are rotated by $-\gamma$ and $+\gamma$ respectively:
\begin{align}
    &\vb{r_{inv,pos,right}}(\phi) = \vb{R}(-\gamma)\;\vb{r_{inv}}(\phi)
\nonumber\\[1ex]
    &\phantom{\vb{r_{inv,po}}}= \frac{d_b}{2} \begin{bmatrix}
        \phantom{-}\cos(\gamma)\;\cos(\phi) +\cos(\gamma)\;\phi \sin(\phi) + \sin(\gamma)\;\sin(\phi) - \sin(\gamma)\;\phi\cos(\phi)\\
        -\sin(\gamma)\;\cos(\phi) - \sin(\gamma)\;\phi \sin(\phi) + \cos(\gamma)\;\sin(\phi) - \cos(\gamma)\;\phi\cos(\phi)\\
    \end{bmatrix}
\\[2ex]
    &\vb{r_{inv,pos,left}}(\phi) = \vb{R}(\gamma)\;\vb{r_{inv}}(\phi)
\nonumber\\[1ex]
    &\phantom{\vb{r_{inv,po}}} = \frac{d_b}{2}
    \begin{bmatrix}
        \phantom{-}\cos(\phantom{-}\gamma)\;\cos(\phi) +\cos(\phantom{-}\gamma)\;\phi \sin(\phi) + \sin(-\gamma)\;\sin(\phi) - \sin(-\gamma)\;\phi\cos(\phi)\\
        -\sin(-\gamma)\;\cos(\phi) - \sin(-\gamma)\;\phi \sin(\phi) + \cos(\phantom{-}\gamma)\;\sin(\phi) - \cos(\phantom{-}\gamma)\;\phi\cos(\phi)\\
    \end{bmatrix}
\end{align}

\paragraph{Equating Curves}
Setting the two curves equal yields two equations, one for each component.

The $x$-component gives:
\begin{equation}
\begin{aligned}
       \cos(\gamma)\;\cos(\phi_r) &+ \cos(\gamma)\;\phi_r\sin(\phi_r)
     + \sin(\gamma)\;\sin(\phi_r)  - \sin(\gamma)\;\phi_r\cos(\phi_r) \\
    &= \cos(\gamma)\;\cos(\phi_l)  + \cos(\gamma)\;\phi_l\sin(\phi_l) 
     - \sin(\gamma)\;\sin(\phi_l)  + \sin(\gamma)\;\phi_l\cos(\phi_l)
\end{aligned}
\end{equation}

The $y$-component gives:
\begin{equation}
\begin{aligned}
     - \sin(\gamma)\;\cos(\phi_r) &- \sin(\gamma)\;\phi_r\sin(\phi_r)
     + \cos(\gamma)\;\sin(\phi_r)  - \cos(\gamma)\;\phi_r\cos(\phi_r) \\
    &= \sin(\gamma)\;\cos(\phi_l)  + \sin(\gamma)\;\phi_l\sin(\phi_l)
     + \cos(\gamma)\;\sin(\phi_l)  - \cos(\gamma)\;\phi_l\cos(\phi_l)
\end{aligned}
\end{equation}

\subparagraph{Symmetry Argument}
Due to the symmetry of the two curves, the ansatz $\phi_r = -\phi_l = \phi$ is proposed. This can be verified by substituting into the $x$-component equation, which yields:
\begin{equation}
\begin{aligned}
    &\phantom{=}\,\;\cos(\gamma)\;\cos(\phi) + \cos(\gamma)\;\phi \sin(\phi) + \sin(\gamma)\;\sin(\phi) - \sin(\gamma)\;\phi\cos(\phi) \\
    &= \cos(\gamma)\;\cos(\phi) + \cos(\gamma)\;\phi \sin(\phi) + \sin(\gamma)\;\sin(\phi) - \sin(\gamma)\;\phi\cos(\phi)
\end{aligned}
\end{equation}
Both sides are identical, confirming the symmetry assumption is valid.

\subparagraph{Deriving the Intersection Condition}
Inserting $\phi_r = -\phi_l = \phi$ into the $y$-component equation:
\begin{equation}
\begin{aligned}
    &\phantom{=}\,\; -\sin(\gamma)\;\cos(\phi) - \sin(\gamma)\;\phi \sin(\phi) + \cos(\gamma)\;\sin(\phi) - \cos(\gamma)\;\phi\cos(\phi)\\
    &=\phantom{-}\sin(\gamma)\;\cos(\phi) + \sin(\gamma)\;\phi \sin(\phi) - \cos(\gamma)\;\sin(\phi) + \cos(\gamma)\;\phi\cos(\phi)
\end{aligned}
\end{equation}

Taking the difference (RHS $-$ LHS $= 0$) and simplifying:

\begin{subequations}
\begin{align}
    2\sin(\gamma)\cos(\phi) + 2\sin(\gamma)\phi\sin(\phi) - 2\cos(\gamma)\sin(\phi) + 2\cos(\gamma)\phi\cos(\phi) &= 0\\
    \sin(\gamma)\cos(\phi) + \sin(\gamma)\phi\sin(\phi) - \cos(\gamma)\sin(\phi) + \cos(\gamma)\phi\cos(\phi) &= 0\\
    \sin(\gamma)\left[\cos(\phi) + \phi\sin(\phi)\right] - \cos(\gamma)\left[\sin(\phi) - \phi\cos(\phi)\right]&= 0
\end{align}
\vspace*{-2.5em}
\begin{align}
    \sin(\gamma)\left[\cos(\phi) + \phi\sin(\phi)\right] &= \cos(\gamma)\left[\sin(\phi) - \phi\cos(\phi)\right]\\
    \tan(\gamma) &= \frac{\sin(\phi) - \phi\cos(\phi)}{\cos(\phi) + \phi\sin(\phi)}\\
    \frac{\sin(\phi) - \phi\cos(\phi)}{\cos(\phi) + \phi\sin(\phi)} - \tan(\gamma) &=0
\end{align}
\end{subequations}

\paragraph{Applying Newton-Raphson}
This means we need to solve the following equation with Newton-Raphson:
\begin{align}
    f(\phi) = \frac{\sin(\phi) - \phi\cos(\phi)}{\cos(\phi) + \phi\sin(\phi)} - \tan(\gamma) &=0
\end{align}
For this we need to compute the derivative of $f(\phi)$ first.

\subparagraph{Derivative of $f(\phi)$}
\begin{enumerate}[label=\alph*)]
    \item Split into components
        \begin{equation}
            \begin{gathered}
                f(\phi) = \frac{u(\phi)}{v(\phi)} -C\\
                \begin{aligned}
                    u(\phi) &=\sin(\phi) - \phi \cos(\phi) \\
                    v(\phi) &=\cos(\phi) + \phi \sin(\phi) \\
                    C &= \tan(\gamma)
                \end{aligned}
            \end{gathered}
        \end{equation}
    \item compute $u'(\phi)$
        \begin{equation}
            \begin{aligned}
                u'(\phi)    &= \cos(\phi) - \cos(\phi) + \phi\sin(\phi) \\
                            &= \phi\sin(\phi)
            \end{aligned}
        \end{equation}
    \item compute $v'(\phi)$
        \begin{equation}
            \begin{aligned}
                v'(\phi)    &= -\sin(\phi) + \sin(\phi) + \phi\cos(\phi) \\
                            &= \phi\cos(\phi)
            \end{aligned}
        \end{equation}
    \item apply quotient rule
        \begin{equation}
            \dv{}{\phi}\left[\frac{u(\phi)}{v(\phi)}\right] = \frac{u'\;v - u\;v'}{v^2}
        \end{equation}
        \begin{subequations}
        \begin{align}
            \dv{f(\phi)}{\phi} & = \frac{u'\;v - u\;v'}{v^2} \\
                &= \frac{
                    \phi\sin(\phi)\cdot\left[\cos(\phi) + \phi\sin(\phi)\right] - \phi\cos(\phi)\cdot\left[\sin(\phi) - \phi\cos(\phi)\right]
                    }
                    {\left[\cos(\phi) + \phi\sin(\phi)\right]^2}\\[1ex]
                &= \frac{
                    \cancel{\phi\sin(\phi)\cos(\phi)} + \phi^2\sin^2(\phi) \cancel{ -\phi\sin(\phi)\cos(\phi)} + \phi^2\cos^2(\phi)
                    }
                    {\left[\cos(\phi) + \phi\sin(\phi)\right]^2}\\[1ex]
                &= \frac{
                    \phi^2\cancelto{1}{\left[\sin^2(\phi) + \cos^2(\phi)\right]}
                    }
                    {\left[\cos(\phi) + \phi\sin(\phi)\right]^2}\\[1ex]
                &= \frac{\phi^2}{\left[\cos(\phi) + \phi\sin(\phi)\right]^2}
        \end{align}
        \end{subequations}
\end{enumerate}

\subparagraph{Newton-Raphson Iteration}
\begin{subequations}
\begin{align}
    \phi_{n+1}  &= \phi_n - \frac{F(\phi_n)}{F'(\phi_n)} \\
                &= \phi_n - \left(\frac{\sin(\phi_n) - \phi_n\cos(\phi_n)}{\cos(\phi_n) + \phi_n\sin(\phi_n)} - \tan(\gamma)\right)
                    \frac{\left[\cos(\phi_n) + \phi_n\sin(\phi_n)\right]^2}{\phi_n^2}\\[1ex]
                &= \phi_n - \frac{\cos(\phi_n) + \phi_n\sin(\phi_n)}{\phi_n^2}\left(\sin(\phi_n) 
                    -\phi_n\cos(\phi_n) - \tan(\gamma)\;\left[\cos(\phi_n) + \phi_n\sin(\phi_n)\right]\right)
\end{align}
\end{subequations}
In a more compact form:
\begin{equation}
    \begin{gathered}
        \phi_{n+1} = \phi_n - \frac{a}{\phi_n^2}\left(b - \tan(\gamma)\;a\right)\\[1ex]
        a =\cos(\phi_n) + \phi_n\sin(\phi_n) \qq{and} b=\sin(\phi_n) - \phi_n\cos(\phi_n)
    \end{gathered}
\end{equation}


\subsubsection{Involute--Undercut Intersection}
In order to properly construct the tooth we need to find the intersection between involute and undercut which usually is just above the pitch circle. Using this we can draw the undercurve upto this intersection and start the involute from this intersection.
\paragraph{System of Equations}
\begin{subequations}
\begin{align}
    \vb{r_{inv}}(\phi_i) =
        \begin{bmatrix}
            x(\phi_i) \\
            y(\phi_i)
        \end{bmatrix}
    &= \frac{d_b}{2}
        \begin{bmatrix}
            \cos(\phi_i) \\
            \sin(\phi_i)
        \end{bmatrix}
        + \frac{d_b}{2}\phi_i
        \begin{bmatrix}
            \sin(\phi_i) \\
            -\cos(\phi_i)
        \end{bmatrix} \\
    \vb{r_{undercut}}(\phi_u) &= \frac{1}{2} \left(
        \begin{bmatrix} a \\ b \end{bmatrix} \cos(\phi_u)
        + \begin{bmatrix} -b \\ a \end{bmatrix}\sin(\phi_u)
        + d_p\,\phi_u\; \begin{bmatrix} \sin\phi_u \\ -\cos\phi_u \end{bmatrix}
    \right)\\[1ex]
    & a = d_f \quad\quad b = \mypm d_f\tan(\alpha_t) \nonumber
\end{align}
\end{subequations}
By looking at how the two curves are defined (see Figures~\ref{fig:involute_construction}~and~\ref{fig:undercut_generated_video}) one can see that the involute needs to be rotated by $\mypm(-\gamma)$ and the undercut curve needs to be rotated by $\mypm(-\gamma - \alpha_t)$. When finding the intersection this is equivalent to rotating the involute by $\mypm \alpha_t$.

Note: the red $\mypm$ refers to either the right flank ($\textcolor{red}{+}$) or left flank ($\textcolor{red}{-}$).

Thus the rotated involute equation is:
\begin{subequations}
\begin{align}
    \vb{r_{inv,pos}}(\phi_i) &= \vb{R}(\mypm\alpha_t)\;\vb{r_{inv}}(\phi_i) \\[1ex]
    &=  \frac{d_b}{2} \begin{bmatrix}
            \cos(\alpha_t) & -\sin(\mypm\alpha_t) \\
            \sin(\mypm\alpha_t) & \cos(\alpha_t)
        \end{bmatrix}
        \begin{bmatrix}
            \cos(\phi_i) + \phi_i\sin(\phi_i)\\
            \sin(\phi_i) -\phi_i \cos(\phi_i)
        \end{bmatrix}
\\
    &=  \frac{d_b}{2}\begin{bmatrix}
            \cos(\alpha_t)\cos(\phi_i) +\cos(\alpha_t)\phi_i\sin(\phi_i) -\sin(\mypm\alpha_t)\sin(\phi_i) +\sin(\mypm\alpha_t)\phi_i\cos(\phi_i) \\
            \sin(\mypm\alpha_t)\cos(\phi_i) +\sin(\mypm\alpha_t)\phi_i\sin(\phi_i) + \cos(\alpha_t)\sin(\phi_i) -\cos(\alpha_t)\phi_i\cos(\phi_i) \\
        \end{bmatrix}
\end{align}
\end{subequations}

\paragraph{Equating Curves}
\begin{subequations}
\begin{equation}
    \vb{r_{inv,pos}}(\phi_i) = \vb{r_{undercut}}(\phi_u)
\end{equation}

\begin{equation}
\begin{aligned}
    &   \frac{d_b}{2}\begin{bmatrix}
            c\cos(\phi_i) + c\phi_i\sin(\phi_i) - d\sin(\phi_i) + d\phi_i\cos(\phi_i) \\
            d\cos(\phi_i) + d\phi_i\sin(\phi_i) + c\sin(\phi_i) - c\phi_i\cos(\phi_i) \\
        \end{bmatrix}\\
    &   \hspace{10em} = \frac{1}{2}\left(\begin{bmatrix}
            a \\ b
        \end{bmatrix}
        \cos\phi_u + \begin{bmatrix}
            -b \\ a
        \end{bmatrix}
        \sin\phi_u + d_p\phi_u\begin{bmatrix}
            \sin\phi_u \\ -
            \cos\phi_u
        \end{bmatrix}\right)
\end{aligned}
\end{equation}

\begin{equation}
    d_b\begin{bmatrix}
        c\cos(\phi_i) + c\phi_i\sin(\phi_i) - d\sin(\phi_i) + d\phi_i\cos(\phi_i) \\
        d\cos(\phi_i) + d\phi_i\sin(\phi_i) + c\sin(\phi_i) - c\phi_i\cos(\phi_i) \\
    \end{bmatrix}
    =
    \begin{bmatrix}
        a \cos(\phi_u) - b\sin(\phi_u) + d_p\phi_u\sin(\phi_u) \\
        b \cos(\phi_u) + a\sin(\phi_u) - d_p\phi_u\cos(\phi_u)
    \end{bmatrix}
\end{equation}
\end{subequations}

\subparagraph{Equating in $\vb x$}
\begin{equation}
\begin{aligned}
    f_1     &= d_bc\cos(\phi_i) + d_bc\phi_i\sin(\phi_i) - d_bd\sin(\phi_i) + d_bd\phi_i\cos(\phi_i) \\
            &\phantom{=}- a\cos(\phi_u) + b\sin(\phi_u) - d_p\phi_u\sin(\phi_u)\\
            &= 0
\end{aligned}
\end{equation}

\subparagraph{Equating in $\vb y$}
\begin{equation}
\begin{aligned}
    f_2     &= d_bd\cos(\phi_i) + d_bd\phi_i\sin(\phi_i) + d_bc\sin(\phi_i) - d_bc\phi_i\cos(\phi_i) \\
            &\phantom{=}- b\cos(\phi_u) - a \sin(\phi_u) + d_p\phi_u\cos(\phi_u) \\
            &= 0
\end{aligned}
\end{equation}

\paragraph{Applying Newton-Raphson}

\subparagraph{Jacobian Matrix}
Before we can numerically solve this using Newton-Raphson we first need to compute the Jacobian of this system of equations:
\begin{align}
    x_1     = \phi_i \qq{,} x_2=\phi_u \\[3ex]
    \vb J   =   \begin{bmatrix}
                    \pdv{f_1}{x_1} & \pdv{f_1}{x_2} \\
                    \pdv{f_2}{x_1} & \pdv{f_2}{x_2}
                \end{bmatrix}
\end{align}

\begin{equation}
\begin{aligned}
    f_1 &= d_bc\cos(x_1) + d_bcx_1\sin(x_1) - d_bd\sin(x_1) + d_bdx_1\cos(x_1) \\
        &\phantom{=}- a\cos(x_2) + b\sin(x_2) - d_px_2\sin(x_2)
\end{aligned}
\end{equation}

\begin{equation}
\begin{aligned}
    f_2 &= d_bd\cos(x_1) + d_bdx_1\sin(x_1) + d_bc\sin(x_1) -d_bcx_1\cos(x_1) \\
        &\phantom{=}-b\cos(x_2) - a \sin(x_2) + d_px_2\cos(x_2)
\end{aligned}
\end{equation}

\begin{subequations}
\begin{align}
    \pdv{f_1}{x_1} &= \cancel{-d_bc\sin(x_1) + d_bc\sin(x_1)} + d_bcx_1\cos(x_1) \cancel{- d_bd\cos(x_1) + d_bd\cos(x_1)} - d_bdx_1\sin(x_1)\\
        &= d_bcx_1\cos(x_1) - d_bdx_1\sin(x_1)\\
        &= d_bx_1[c\cos(x_1) - d\sin(x_1)]\\
            \pdv{f_1}{x_2} &= a\sin x_2 + b\cos x_2 - d_p \sin x_2 - d_px_2 \cos x_2 \\
            \pdv{f_2}{x_1} &= \cancel{-d_bd\sin(x_1) + d_bd\sin(x_1)} + d_bdx_1\cos(x_1) \cancel{+ db_c\cos(x_1) - d_bc\cos(x_1)} + d_bcx_1\sin(x_1)\\
        &= d_bdx_1\cos(x_1) + d_bcx_1\sin(x_1)\\
        &= d_bx_1[d\cos(x_1) + c\sin(x_1)]\\
            \pdv{f_2}{x_2} &= b\sin x_2 - a\cos x_2 + d_p\cos x_2 - d_px_2\sin x_2
\end{align}
\end{subequations}
Finally:
\begin{align}
    \vb J =
        \begin{bmatrix}
            d_bx_1[c\cos x_1 - d\sin x_1] & a\sin x_2 + b\cos x_2 - d_p \sin x_2 - d_px_2 \cos x_2\\
            d_bx_1[d\cos x_1 + c\sin x_1] & b\sin x_2+a\cos x_2 + d_p\cos x_2 - d_px_2\sin x_2 \\
        \end{bmatrix}
\end{align}

\subparagraph{Inverse of Jacobian}
We compute the inverse of the Jacobian as:
\begin{align}
    \vb {J^{-1}} &= \frac{1}{J_{11}J_{22} - J_{12}J_{21}}
        \begin{bmatrix}
            J_{22} & - J_{12} \\
            -J_{21} & J_{11}
        \end{bmatrix}\\[2ex]
    J_{11} &= d_bx_1[c\cos x_1 - d\sin x_1] \\
    J_{12} &= a\sin x_2 + b\cos x_2 - d_p \sin x_2 - d_px_2 \cos x_2  \\
    J_{21} &= d_bx_1[d\cos x_1 + c\sin x_1] \\
    J_{22} &= b\sin x_2 - a\cos x_2 + d_p\cos x_2 - d_px_2\sin x_2
\end{align}

\subparagraph{Newton-Raphson Iteration}
Finally the intersection can be found by iterating:
\begin{align}
    \begin{bmatrix}
        x_1 \\ x_2
    \end{bmatrix}^{(k+1)}
    =
    \begin{bmatrix}
        x_1 \\ x_2
    \end{bmatrix}^{(k)}
    - \vb{J^{-1}} \cdot \begin{bmatrix} f_1 \\ f_2 \end{bmatrix}
\end{align}
where
\begin{align}
    x_1 = \phi_i \qq{,} x_2=\phi_u \\[3ex]
    f_1 &= d_bc\cos(x_1) + d_bcx_1\sin(x_1) - d_bd\sin(x_1) + d_bdx_1\cos(x_1)\nonumber\\
        &\phantom{=}- a\cos(x_2) + b\sin(x_2) - d_px_2\sin(x_2)\\[1ex]
    f_2 &= d_bd\cos(x_1) + d_bdx_1\sin(x_1) + d_bc\sin(x_1) -d_bcx_1\cos(x_1)\nonumber\\
        &\phantom{=}-b\cos(x_2) - a \sin(x_2) + d_px_2\cos(x_2) \\[3ex]
    \vb {J^{-1}} &= \frac{1}{J_{11}J_{22} - J_{12}J_{21}}
        \begin{bmatrix}
            J_{22} & - J_{12} \\
            -J_{21} & J_{11}
        \end{bmatrix}\\[2ex]
    J_{11} &= d_bx_1[c\cos x_1 - d\sin x_1] \\
    J_{12} &= a\sin x_2 + b\cos x_2 - d_p \sin x_2 - d_px_2 \cos x_2  \\
    J_{21} &= d_bx_1[d\cos x_1 + c\sin x_1] \\
    J_{22} &= b\sin x_2 - a\cos x_2 + d_p\cos x_2 - d_px_2\sin x_2
\end{align}

\subsection{Summary}
\label{sec:spur_summary}
This section consolidates all final formulas needed to construct a single spur gear tooth. For derivations, refer to the corresponding sections.

\subsubsection{Input Parameters and Derived Quantities}

\begin{center}
\begin{tabular}{@{}llll@{}}
    \textbf{Symbol} & \textbf{Parameter} & \textbf{Default} & \textbf{Unit} \\
    \hline
    $m$         & Module                    & --            & mm \\
    $z$         & Number of teeth           & --            & -- \\
    $\alpha_n$  & Normal pressure angle     & $\SI{20}{\degree}$ & deg \\
    $\beta$     & Helix angle               & $\SI{0}{\degree}$ (spur) & deg \\
    $x$         & Profile shift coefficient & 0             & -- \\
    $h_a^*$     & Addendum coefficient      & 1             & -- \\
    $c^*$       & Tip clearance factor      & 0.25          & -- \\
\end{tabular}
\end{center}

From these, compute the circle diameters (Section~\ref{sec:preliminaries}) and half tooth angle (Section~\ref{subsec:tooth_thickness}):
\begin{align}
    \alpha_t &= \arctan\!\left(\frac{\tan(\alpha_n)}{\cos(\beta)}\right) \\
    d_p &= m \cdot z \\
    d_b &= d_p \cos(\alpha_t) \\
    d_a &= m\cdot(z + 2\,x + 2\cdot h_a^*) \\
    d_f &= m\cdot(z + 2\,x - 2\cdot(h_a^*+c^*)) \\
    \gamma &= \frac{m\left(\pi + 2 \cdot x \cdot \tan(\alpha_n)\right)}{2\,d_p} + \left[\left(\frac{d_p}{d_b}\right)^2 - 1\right]^{\tfrac{1}{2}} - \arctan\left(\left[\left(\frac{d_p}{d_b}\right)^2 - 1\right]^{\tfrac{1}{2}}\right)
\end{align}

\subsubsection{Parametric Curves}

\paragraph{Involute} (Section~\ref{sec:involute}):
\begin{equation}
    \vb{r_{inv}}(\phi) = \frac{d_b}{2} \begin{bmatrix}
        \cos(\phi) + \phi\sin(\phi) \\
        \sin(\phi) - \phi\cos(\phi)
    \end{bmatrix}
    \qq{,}
    \phi(d^*) = \mypm\left[\left(\frac{d^*}{d_b}\right)^2 - 1\right]^{\frac{1}{2}}
\end{equation}
Start: $\phi_{start} = 0$ (base circle). \quad End: $\phi_{end} = \phi(d_a)$ (addendum circle).

\paragraph{Positioned involute} (Section~\ref{sec:tooth_geometry_positioning}, rotated by $\mypm(-\gamma)$):
\begin{equation}
    \vb{r_{inv,pos}}(\phi) = \frac{d_b}{2} \begin{bmatrix}
        \phantom{-}\cos(\myppm\gamma)\;\cos(\phi) + \cos(\myppm\gamma)\;\phi\sin(\phi)
                 + \sin(\mypm \gamma)\;\sin(\phi) - \sin(\mypm \gamma)\;\phi\cos(\phi) \\
                 - \sin(\mypm \gamma)\;\cos(\phi) - \sin(\mypm \gamma)\;\phi\sin(\phi)
                 + \cos(\myppm\gamma)\;\sin(\phi) - \cos(\myppm\gamma)\;\phi\cos(\phi) \\
    \end{bmatrix}
\end{equation}

\paragraph{Undercut curve} (Section~\ref{sec:undercut}):
\begin{equation}
    \begin{gathered}
        \vb{r_{undercut}}(\phi) = \frac{1}{2}\left(
              \begin{bmatrix}  a \\ b \end{bmatrix}\cos(\phi)
            + \begin{bmatrix} -b \\ a \end{bmatrix}\sin(\phi)
            + d_p\,\phi\;\begin{bmatrix} \sin(\phi) \\ -\cos(\phi) \end{bmatrix}
        \right)
    \\
        a = d_f \quad\quad b = \mypm d_f\tan(\alpha_t)
    \end{gathered}
\end{equation}
Start: $\phi_0 = \mypm\frac{d_f}{d_p}\tan(\alpha_t)$ (dedendum circle).

\paragraph{Positioned undercut} (Section~\ref{sec:tooth_geometry_positioning}, rotated by $\mypm(-\gamma - \alpha_t)$):
\begin{subequations}
\begin{align}
    &\vb{r_{undercut,pos}}(\phi) = \vb{R}(\mypm(-\gamma - \alpha_t))\;\vb{r_{undercut}}(\phi)
\\[1ex]
    &= \frac{1}{2}\bigg[\begin{array}{cc}
            \phantom{-}a\cos(\myppm\;\gamma+\alpha_t\;)\;\cos(\phi)
            - b\cos(\myppm\;\gamma+\alpha_t\;)\;\sin(\phi)
                + d_p\cos(\myppm\;\gamma+\alpha_t\;)\;\phi\sin(\phi)\dots \\
            -a\sin(\mypm(\gamma+\alpha_t))\;\cos(\phi)
                +b\sin(\mypm(\gamma+\alpha_t))\;\sin(\phi)
                - d_p\sin(\mypm(\gamma+\alpha_t))\;\phi\sin(\phi)\dots
        \end{array}
\nonumber\\
        &\phantom{= \frac{1}{2}\bigg[}\begin{array}{cc}
            \hspace{1em}\dots + b\sin(\mypm(\gamma+\alpha_t))\;\cos(\phi)
                + a\sin(\mypm(\gamma+\alpha_t))\;\sin(\phi)
                - d_p\sin(\mypm(\gamma+\alpha_t))\;\phi\cos(\phi)\\
            \hspace{1em}\dots+ b\cos(\myppm\;\gamma+\alpha_t\;)\;\cos(\phi)
                + a\cos(\myppm\;\gamma+\alpha_t\;)\;\sin(\phi)
                - d_p\cos(\myppm\;\gamma+\alpha_t\;)\;\phi\cos(\phi)
        \end{array}\bigg]
\end{align}
\end{subequations}

\subsubsection{Curve Intersections}

\paragraph{Left--Right Involute Intersection} (Section~\ref{sec:curve_intersections}). Solve with 1D Newton-Raphson (Eq.~\eqref{eq:1D_newton_raphson}):
\begin{equation}
    \begin{gathered}
        \phi_{n+1} = \phi_n - \frac{a}{\phi_n^2}\left(b - \tan(\gamma)\;a\right) \\[1ex]
        a = \cos(\phi_n) + \phi_n\sin(\phi_n) \qq{,} b = \sin(\phi_n) - \phi_n\cos(\phi_n)
    \end{gathered}
\end{equation}
If the intersection happens before the dedendum circle ($\phi_\infty < \phi(d_a)$), the involute end is updated to $\phi_{end} = \phi_\infty$ instead of $\phi(d_a)$.

\paragraph{Involute--Undercut Intersection} Solve with 2D Newton-Raphson (Eq.~\eqref{eq:2D_newton_raphson}), using $x_1 = \phi_i$ (involute), $x_2 = \phi_u$ (undercut), $c = \cos(\alpha_t)$, $d = \sin(\mypm\alpha_t)$:
\begin{equation}
    \begin{bmatrix} x_1 \\ x_2 \end{bmatrix}^{(k+1)}
    = \begin{bmatrix} x_1 \\ x_2 \end{bmatrix}^{(k)}
    - \vb{J^{-1}} \cdot \begin{bmatrix} f_1 \\ f_2 \end{bmatrix}
\end{equation}
\begin{align}
    f_1 &= d_bc\cos(x_1) + d_bcx_1\sin(x_1) - d_bd\sin(x_1) + d_bdx_1\cos(x_1) \nonumber\\
        &\phantom{=}- a\cos(x_2) + b\sin(x_2) - d_px_2\sin(x_2) \\
    f_2 &= d_bd\cos(x_1) + d_bdx_1\sin(x_1) + d_bc\sin(x_1) - d_bcx_1\cos(x_1) \nonumber\\
        &\phantom{=}- b\cos(x_2) - a\sin(x_2) + d_px_2\cos(x_2)
\end{align}
\begin{equation}
    \vb{J^{-1}} = \frac{1}{J_{11}J_{22} - J_{12}J_{21}}
    \begin{bmatrix}
        J_{22} & -J_{12} \\
        -J_{21} & J_{11}
    \end{bmatrix}
\end{equation}
\begin{align}
    J_{11} &= d_bx_1[c\cos x_1 - d\sin x_1] &
    J_{12} &= a\sin x_2 + b\cos x_2 - d_p\sin x_2 - d_px_2\cos x_2 \nonumber\\
    J_{21} &= d_bx_1[d\cos x_1 + c\sin x_1] &
    J_{22} &= b\sin x_2 - a\cos x_2 + d_p\cos x_2 - d_px_2\sin x_2
\end{align}
The intersection determines the actual start of the involute ($\phi_i$) and the end of the undercut ($\phi_u$).
